\documentclass[noauthor,noinstructornotes]{ximera}

\title{Week 6}
\author{Math 1193 Design Team}

\begin{document}
\begin{abstract}
    Activities for Week 6.
\end{abstract}
\maketitle

\begin{problem}
Look over the following table classifying expressions.

\begin{center}
\begin{tabular}{|c|c|c|c|} 
 \hline
 Expression & Polynomial? & Degree & Leading Coefficient \\
 \hline\hline
 \(x\) & Yes & 1 & 1 \\ 
 \hline
 \(3x^2\) & Yes & 2 & 3 \\
 \hline
 \(7x + x^{1/2}\) & No & N/A & N/A \\
 \hline
 \(2x^3 - x^2 + x - 4\) & Yes & 3 & 2 \\
 \hline
 \(\pi x^2 - 9\) & Yes & 2 & \(\pi\) \\
 \hline
 \(\frac{2}{3}x^5 + \frac{9}{7}x\) & Yes & 5 & \(\frac{2}{3}\) \\
 \hline
 \(\frac{2x^5 + 9x}{3x^5 - 7x}\) & No & N/A & N/A \\
 \hline
 \(x^\pi - 2\) & No & N/A & N/A \\
 \hline
 \(2^x + 2^x + 1\) & No & N/A & N/A \\
 \hline
 \((2 - x)(x + 1)\) & Yes & 2 & \(-1\) \\
 \hline
\end{tabular}
\end{center}

\begin{enumerate}
    \item Come up with a definition of a polynomial. Reviewing the
    definition of a quadratic may help you. Check with your group
    and the instructor to confirm.
    \item Identify all of the following which are polynomials 
    and give their degree and leading coefficient.
    \begin{enumerate}
        \item \(f\left(x\right)=2x^3-x+1\)
	    \item \(g\left(x\right)=2x^\pi-x+1\)
	    \item \(h\left(x\right)=7-x\)
	    \item \(f\left(x\right)=x\)
	    \item \(g\left(x\right)=\left(x-17\right)\left(x+8\right)\)
	    \item \(h\left(x\right)=\left(2x-3\right)^5\left(x-1\right)^2\)
	    \item \(f\left(x\right)= e x^5-3\)
	    \item \(g\left(x\right)=\frac{x-1}{x+1}\)
	    \item \(h\left(x\right)=7x^6-2x+x^6\)
	    \item \(f\left(x\right)=1+x+x^2\)
	    \item \(g\left(x\right)=-3\left(x-\frac{1}{2}\right)^2+4\)
	    \item \(h\left(x\right)=\left(x^2+1\right)\left(x-1\right)\left(x+1\right)^2\)
	    \item \(f\left(x\right)=\left(-3x+2\right)\left(5x+1\right)\left(2x-2\right)^2\) 
    \end{enumerate}
\end{enumerate}
\end{problem}

\begin{problem}
Solve the following polynomial equations.

\begin{enumerate}
    \item \((x - 1)(x + 2) = 0\)
    \item \((x - 1)^2(x + 2) = 0\)
    \item \(-3(x - 1)^2(x + 2)\)
    \item \((2x - 1)^3(x + 1)\)
    \item \(2x^3 + 8x^2 + 3x + 12 = 0\)
    \item \(2x^2 + 7x + 3 = 0\)
\end{enumerate}

How could you check your answers? 

Which equations were most difficult to solve?
\end{problem}

\begin{problem}
Ella is working on some math homework and decides to graph the 
function 
\(f\left(x\right)=\left(x-2\right)\left(x+3\right)\left(x-5\right)\left(x+15\right)\) 
on her calculator to get a sense of what the graph looks like. Her 
viewing window on the graphing calculator is set to \([-10,10]\) by \([-10,10]\). 

\begin{enumerate}
    \item How many \(x\)-intercepts will Ella see with this viewing window?
    \item How many \(x\)-intercepts does \(f\) have?
    \item What viewing window would you use?
\end{enumerate}

\begin{instructorNotes}
This was inspired by the textbook and highlights a pitfall when relying 
too heavily on the calculator. 
\end{instructorNotes}
\end{problem}

\begin{problem}
Consider the following graph of the function \(f\).

\begin{center}
    \includegraphics[width=10cm]{endBehaviorLinear.png}
\end{center}

\begin{enumerate}
    \item Write a sentence describing what happens to \(f\left(x\right)\) 
    as \(x\) gets bigger and bigger in the positive direction. 
    Remember that \(f(x)\) is the \(y\)-value associated with \(x\). 
    An example sentence might be ``As \(x\) gets bigger and bigger in 
    the positive direction, \(f(x)\) gets...''.
    \item Write a sentence describing what happens to \(f\left(x\right)\) 
    as \(x\) gets bigger and bigger in the negative direction. 
\end{enumerate}

\end{problem}

\begin{problem}
Consider the following graph of the function \(f\).

\begin{center}
    \includegraphics[width=10cm]{endBehaviorQuadratic.png}
\end{center}

\begin{enumerate}
    \item Write a sentence describing what happens to \(f\left(x\right)\) 
    as \(x\) gets bigger and bigger in the positive direction. 
    Remember that \(f(x)\) is the \(y\)-value associated with \(x\). 
    An example sentence might be ``As \(x\) gets bigger and bigger in 
    the positive direction, \(f(x)\) gets...''.
    \item Write a sentence describing what happens to \(f\left(x\right)\) 
    as \(x\) gets bigger and bigger in the negative direction. 
\end{enumerate}

\end{problem}

\begin{problem}
Consider the polynomial \(g\left(x\right)=2x^5-5x^4-30x^3+5x^2+88x+60\).

\begin{enumerate}
    \item Identify the degree of the polynomial.
    \item Which of the 6 terms, \(2x^5\), \(5x^4\), \(30x^3\), \(5x^2\),
    \(88x\), or \(60\), is greatest when:
    \begin{enumerate}
        \item \(x = 0\)
        \item \(x = 1\)
        \item \(x = 3\)
        \item \(x = 5\)
    \end{enumerate}
    \item Write a sentence describing what happens to \(g\left(x\right)\) 
    as \(x\) gets bigger and bigger in the positive direction. 
    Remember that \(g(x)\) is the \(y\)-value associated with \(x\). 
    An example sentence might be ``As \(x\) gets bigger and bigger in 
    the positive direction, \(g(x)\) gets...''.
    \item Write a sentence describing what happens to \(g\left(x\right)\) 
    as \(x\) gets bigger and bigger in the negative direction. 
\end{enumerate}
\end{problem}

\end{document}
